\section{A blocker moves when its holder is short of cash}
\label{sec:liquidity}

\subsection{Refusal falls with the holder's cash}

Figure~\ref{fig:liquidity}a splits the 289 completion requests by the holder's cash when the request arrives. Holders with under \$100 reject 25\% of them. Holders with \$100 to \$299 reject 49\%, and the two richer groups reject 70\% and 68\%. Counters make up most of the difference. Acceptance on the spot stays rare at every level, at 12\% for the poorest holders and 3\% for the richest.

\begin{figure}[t]
\centering\includegraphics[width=\linewidth]{fig_liquidity.pdf}
\caption{Liquidity and release. (a)~Responses to completion requests by the holder's cash. (b)~Holder's cash when it first rejected or countered the buyer's request and when it later sold, for the one-sided sales with such an earlier refusal. Each line is one holder, buyer and color group. The black line joins the medians.}
\label{fig:liquidity}
\end{figure}

The contrast between holders under \$100 and holders with \$300 or more is 25\% against 69\%. It has the same sign in each of the five seed blocks and in 13 of the 16 games that contain both kinds of request. With any one block removed it ranges from 39 to 49 percentage points.

The contrast shrinks under composition checks and does not vanish. Within opening offers the two shares are 50\% and 80\%, and within later offers 9\% and 30\%. Without Opus as holder they are 35\% and 79\%, and without Opus on either side 54\% and 79\%. Without the offers that follow a conflicting Grok call they are 23\% and 69\%. Sol and Grok show the pattern individually. Sol rejects 15\% of requests when it holds under \$100 and 56\% when it holds \$300 or more, and Grok 27\% and 88\%. Gemini is the exception. It rejects 10 of 11 requests even when it holds under \$100. Opus rejects almost nothing at any level.

Low cash also softens refusals of deeds that complete nothing, where the shares are 23\% and 49\%. Part of the association is therefore a general effect of needing money. Buyers also send completion requests where cash is short. A holder with under \$100 receives 24\% of them, against 8\% of requests for deeds that complete nothing.

\subsection{Following each pair over time}

An offer is a poor unit for a claim about release, because one stubborn pair can produce twenty of them. We therefore group offers into pairs. A pair is a game, a would-be completer, a holder and a color group, with at least one legal offer, written by either party, under which the completer would obtain the group from the holder. The cohort has 148 pairs. A trade settles 63 of them, which are the 63 groups completed by trade. In 43 the sale completes a group for the buyer only, and we call these one-sided sales. In 20 a swap completes a group for both parties.

One-sided sales come disproportionately from holders short of cash. In 20 of the 43 the holder has under \$100 at settlement, and the median is \$178. In the 20 swap completions the median is \$594.50, and 1 holder has under \$100.

The sharper evidence follows a pair through time. In 14 one-sided sales the holder had rejected or countered a request from the same buyer on an earlier turn. Figure~\ref{fig:liquidity}b shows the holder's cash at that first refusal and at the sale, a median of 17.5 turns later. It is lower at the sale in 13 of the 14 pairs and higher in 1, and the median falls from \$571.50 to \$82. The comparison holds the holder, the buyer and the group fixed. All four systems appear as holders, and with Opus's 3 pairs removed the count is 11 of 11. The pairs come from 10 games in 4 of the five seed blocks (Table~\ref{tab:dyads}). Four more one-sided sales settled on a later turn after the buyer, not the holder, had turned down earlier terms. Counting all 18 from the pair's first offer, the holder's cash is lower at the sale in 14. In the 38 pairs that drew offers on more than one turn and never settled, the holder's cash moved much less, from a median of \$671.50 at the first offer to \$516 at the last.

This is not a controlled contrast. Turns pass between the two moments, the board changes, and holders who ran short may differ from those who did not. It is the pattern expected if a blocker is a reserve that its holder gives up when cash runs out, and the agents describe it that way. Sol's notes say it will revisit monopoly trades when opponents need cash (Incident~\ref{inc:g15}) and that it would \InlineQuote{in-g1-d313-distress} (G1 d313).

\subsection{Short of cash is not the same as cheap}

Among one-sided sales for cash alone, the 8 made by Opus with under \$100 settled at a median of 0.66 times list price. The 6 made by the other systems with under \$100 settled at a median of 2.25 times list. Of Opus's sales, 7 were below list price, and of the others' sales 5 were above it. Opus made 12 of the 20 low-cash one-sided sales, nearly all under action-only recovery. As a holder it settled 24 of its 35 pairs, against 13 of 34 for Sol, 13 of 38 for Gemini and 13 of 41 for Grok. The cohort cannot say how much of this belongs to the model and how much to its response condition. It does show that cash pressure and a low price are separable. The other three systems released deeds when short of cash and still charged for the buyer's monopoly.

\IncidentTitle{g1}{A completing deed sold for a quarter of its list price (G1)}

Opus holds Boardwalk and Grok holds Park Place. From G1 turn 69, when it has \$453, Opus offers Boardwalk seven times in exchange for New York Avenue, a swap that would complete orange for Opus as well as dark blue for Grok. Grok rejects six of the offers and counters one. At turn 122 Opus pays a \$550 rent by mortgaging a railroad, which leaves it \$18, and offers the mortgaged deed for cash alone.
\TranscriptQuote{g1-d306-offer}
\TranscriptQuote{g1-d307-counter}
Opus asks \$100 and adds that \InlineQuote{in-g1-d308-weapon}. Grok accepts above its own ceiling.
\TranscriptQuote{g1-d309-accept-note}
\IncidentOutcome{Grok completes dark blue for \$100, a quarter of Boardwalk's list price. It never redeems the mortgage or builds. Over the next 94 turns opponents land on the group five times and pay \$70 in total. Opus is eliminated four turns after the sale, owing Gemini \$550. Gemini wins.}
\IncidentReading{Between its first offer and the sale the holder's cash fell from \$453 to \$18, and its terms fell from a deed that would complete its own group to \$100. Opus presents the sale as arming a rival of the leader. Grok's note prices the alternative, in which Opus goes bankrupt and Boardwalk passes to the leader with the estate. Opus's side has no private note, because it played under action-only recovery.}
\IncidentSource{G1 d305--d315. Grok's responses are standard first attempts. Opus's are action-only responses. The quoted public messages appear in the counterparty's next prompt.}

\IncidentTitle{g13}{Distress used as a bargaining position (G13)}

At G13 turn 107 Grok has \$35, owns the three yellows and stands nine squares before Short Line, where Sol's four railroads charge \$200. Its Ventnor Avenue is unmortgaged and would raise \$130. Grok offers Opus its mortgaged Pennsylvania Avenue for \$200.
\TranscriptQuote{g13-d250-insurance}
Opus counters at \$130. Grok holds at \$200 and tells Opus that \$130 would leave it \$35 short. Its private note and the explanation returned with the same response say more.
\TranscriptQuote{g13-d252-floor}
\TranscriptQuote{g13-d252-omission}
\TranscriptQuote{g13-d253-accept}
\IncidentOutcome{Opus pays \$200 and \$16 of mortgage-transfer interest. Seven turns later it buys North Carolina at auction for \$650 and completes green. Grok's next roll misses Short Line. It pays \$200 on another of Sol's railroads at turn 146 and goes bankrupt on Short Line at turn 150, and the yellows pass to Sol. Opus wins.}
\IncidentReading{Grok presents its own distress as the buyer's problem, since Sol would inherit the yellows if Grok failed. It asks \$200 while its note sets a floor of \$165, and the returned explanation records that it left out the Ventnor mortgage. Under this prompt that is an authorized public bluff. Opus pays the asking price, and nothing in the record shows that it believed the claimed limit or that the omission produced the extra \$70. The pitch also overstates the risk. Short Line is nine squares away, a 4-in-36 chance on one roll, as Grok's own note says.}
\IncidentSource{G13 d250--d253. Grok's responses are standard first attempts. Opus's are action-only responses. The returned explanation is the provider's reasoning field, a channel separate from the private note.}
